\documentclass{article}
\usepackage{graphicx}
\usepackage{booktabs}
\usepackage{array}
\usepackage[margin=1in]{geometry}

\begin{document}

\title{GELANG}
\author{}
\date{}
\maketitle

\section*{Imperative Language Specification}

The language grammar is provided in \texttt{grammar.txt} file. Moreover:

\begin{enumerate}
    \item Arithmetic operations are performed on 64bit integer numbers. 
    \begin{itemize}
        \item Division by zero yield a result of 0 and a remainder of 0 as well.
        \item The remainder from division has the same sign as the divisor, just as in Python.
    \end{itemize}
    
    \item Arrays can be declared with integer indices. 
    \begin{itemize}
        \item An identifier \texttt{t[i]} refers to the $i$-th element of the array \texttt{t}
        \item A declaration where the first number is greater than the second is flagged as an error.
        \item A declaration of the form \texttt{t[-10:10]} means declaring an array with 21 elements indexed from -10 to 10
    \end{itemize}
    
    \item Procedures cannot contain recursive calls. Formal parameters are passed by reference (IN-OUT parameters).
    \begin{itemize}
        \item Variables used within a procedure must be either formal parameters or declared inside the procedure.
        \item The name of the array in formal parameters should be preceded by the letter T.
        \item Only previously defined procedures can be called.
        \item The formal parameters of the calling procedure and its local variables can be provided as arguments to the called procedures.
        \item FOR loop iterators cannot be passed as arguments to a procedure.
    \end{itemize}

    \item A FOR loop has a local iterator, which takes values from the value after \texttt{FROM} to the value after \texttt{TO/DOWNTO} in increments of +1 or -1 accordingly.
    \begin{itemize}
        \item The number of iterations of the FOR loop is determined at the beginning and cannot change during execution.
        \item The iterator of the FOR loop cannot be modified within the loop.Trying to do so will result in compilation error.
    \end{itemize}

    \item A REPEAT-UNTIL loop finishes when the condition written after \texttt{UNTIL} is satisfied. The loop will execute at least once.

    \item The \texttt{READ} instruction reads a value from the input and assigns it to a variable, while the \texttt{WRITE} instruction outputs the value of a variable or number to the output.

    \item The remaining instructions are consistent with their meaning in most programming languages.

    \item A \texttt{pidentifier} is described by the regular expression \texttt{[\_a-z]+}.
    
    \item A \texttt{num} is an 64bit integer in decimal notation. 

    \item Lowercase and uppercase letters are distinguished.

    \item The program may include comments starting with \# and valid until the end of the line.
\end{enumerate}

\section*{Virtual Machine}

The virtual machine consists of a program counter $k$ and a sequence of memory cells $p_i$, for $i = 0, 1, 2, \dots$ (for technical reasons, $i \leq 2^{62}$). The memory cell $p_0$ serves as an accumulator. The machine operates on 64bit integer numbers. The machine program consists of a sequence of instructions, which are implicitly numbered starting from zero. At each step, the instruction at index $k$ is executed until a HALT instruction is encountered. Initially, the contents of memory cells are undefined, and the program counter $k$ is set to 0.

Table \ref{tab:instructions} lists the available instructions along with their interpretations and execution costs. Comments can be included in the program using the format \#comment. Whitespace in the code is ignored. Jumping to a nonexistent instruction is treated as an error.

\begin{table}[h]
\centering
\renewcommand{\arraystretch}{1.4}
\begin{tabular}{|c|>{\raggedright\arraybackslash}m{9cm}|c|}
\hline
\textbf{Instruction} & \textbf{Interpretation} & \textbf{Time} \\
\hline
GET $i$  & Reads a number and stores it in memory cell $p_i$, then $k \gets k + 1$ & 100 \\
PUT $i$  & Displays the contents of memory cell $p_i$, then $k \gets k + 1$ & 100 \\
\hline
LOAD $i$  & $p_0 \gets p_i$, then $k \gets k + 1$ & 10 \\
STORE $i$  & $p_i \gets p_0$, then $k \gets k + 1$ & 10 \\
LOADI $i$  & $p_0 \gets p_{p_i}$, then $k \gets k + 1$ & 20 \\
STOREI $i$  & $p_{p_i} \gets p_0$, then $k \gets k + 1$ & 20 \\
\hline
ADD $i$  & $p_0 \gets p_0 + p_i$, then $k \gets k + 1$ & 10 \\
SUB $i$  & $p_0 \gets p_0 - p_i$, then $k \gets k + 1$ & 10 \\
ADDI $i$  & $p_0 \gets p_0 + p_{p_i}$, then $k \gets k + 1$ & 20 \\
SUBI $i$  & $p_0 \gets p_0 - p_{p_i}$, then $k \gets k + 1$ & 20 \\
\hline
SET $x$  & $p_0 \gets x$, then $k \gets k + 1$ & 50 \\
HALF  & $p_0 \gets \lfloor p_0 / 2 \rfloor$, then $k \gets k + 1$ & 5 \\
\hline
JUMP $j$  & $k \gets k + j$ & 1 \\
JPOS $j$  & If $p_0 > 0$, then $k \gets k + j$, otherwise $k \gets k + 1$ & 1 \\
JZERO $j$  & If $p_0 = 0$, then $k \gets k + j$, otherwise $k \gets k + 1$ & 1 \\
JNEG $j$  & If $p_0 < 0$, then $k \gets k + j$, otherwise $k \gets k + 1$ & 1 \\
\hline
RTRN $i$  & $k \gets p_i$ & 10 \\
HALT  & Stops the program & 0 \\
\hline
\end{tabular}
\caption{Instruction set of the virtual machine}
\label{tab:instructions}
\end{table}

\end{document}
